%%%%%%%%%%%%%%%%%%%%%%%%%%%%%%%%%%%%%%%%%%%%%%%%%%%%%%%%%%%%%%%%%%%%%%%%
%
% Two-page extended abstract
% From Human Videos to Dexterous Robot Capability workshop (CoRL 2026)
% Formatting: NVIDIA tech report.
%
% Limit: 2 pages of main text; references, acknowledgments and
% appendices do not count.
%
% Build:  pdflatex main && bibtex main && pdflatex main && pdflatex main
%
%%%%%%%%%%%%%%%%%%%%%%%%%%%%%%%%%%%%%%%%%%%%%%%%%%%%%%%%%%%%%%%%%%%%%%%%

\pdfoutput=1

%%% == IMPORTANT ==
%%% For anonymized submission (no authors, logo, copyright or
%%% acknowledgments), use this:
%\documentclass[10pt, logo, copyright, anonymous]{nvidiaextabstract}

%%% For camera-ready, use this:
\documentclass[10pt]{nvidiaextabstract}

%%% Load required packages here (many are loaded by the class already:
%%% amsmath, amssymb, graphicx, booktabs, xcolor, hyperref, natbib,
%%% balance, enumitem, caption).
\usepackage{microtype}
\usepackage{multirow}
\usepackage{subcaption}
\usepackage{xspace}

%%%%%%%%%%%%%%%%%%%%%%%%%%%%%%%%%%%%%%%%%%%%%%%%%%%%%%%%%%%%%%%%%%%%%%%%

%%% Venue: printed in the first-page footer. Comment out to drop it.
\venue{From Human Videos to Dexterous Robot Capability (CoRL 2026)}

%%% Logo in the top-left corner of page 1 (shown with the `logo' class
%%% option). Options in [] go to \includegraphics; default is the NVIDIA
%%% logo at width=90pt. The header is about 40pt tall, so keep height below that.
%\logo[height=28pt]{assets/mylogo.png}

%%% Printed instead of the authors when the `anonymous' option is set.
\submissionID{<<Registration team name>>}

%%% Title. The subtitle defaults to "Extended Abstract".
\title{Video to Data Challenge Technical Report Template}
%\subtitle{Extended Abstract}
%\shorttitle{Short Running Title}   % running head on page 2

%%% Authors (authblk syntax).
\author[1]{First Author}
\author[1]{Second Author}
\author[2]{Third Author}
\affil[1]{NVIDIA}
\affil[2]{University Name}
\affil[ ]{\texttt{\{first,second\}@nvidia.com}, \texttt{third@university.edu}}

%%% Abstract. Must be in the preamble and must not contain any
%%% \begin...\end environments.
\begin{abstract}
\textbf{Abstract:}
This extended abstract should be two pages, not including references, acknowledgments, and appendices.
NVIDIA logos are placeholders.
\end{abstract}

%%%%%%%%%%%%%%%%%%%%%%%%%%%%%%%%%%%%%%%%%%%%%%%%%%%%%%%%%%%%%%%%%%%%%%%%

%%% Author-defined commands.
\definecolor{NvidiaGreen}{RGB}{118, 185, 0}
\newcommand{\method}{\textsc{Method}\xspace}

%%%%%%%%%%%%%%%%%%%%%%%%%%%%%%%%%%%%%%%%%%%%%%%%%%%%%%%%%%%%%%%%%%%%%%%%

\begin{document}

\maketitle

%%%%%%%%%%%%%%%%%%%%%%%%%%%%%%%%%%%%%%%%%%%%%%%%%%%%%%%%%%%%%%%%%%%%%%%%

\section{Introduction}
\label{sec:intro}

Motivate the problem and place it in the video to data challenge and the
tracks your team participates in. Briefly position your work against the
closest prior work in lieu of a full related works. Cite with
\verb|\cite|, e.g.~\cite{Lam94}.

Close the introduction with the contributions and findings of this work:
\begin{itemize}
  \item the first contribution;
  \item the second contribution;
  \item ...
\end{itemize}

%%%%%%%%%%%%%%%%%%%%%%%%%%%%%%%%%%%%%%%%%%%%%%%%%%%%%%%%%%%%%%%%%%%%%%%%

\section{Approach}
\label{sec:approach}

Introduce the setting and notation. For example, a policy
$\pi_\theta(a_t \mid o_t)$ maps observations $o_t$ to robot actions $a_t$.
Display only the equations you refer back to:
%
\begin{equation}
\label{eq:objective}
J(\theta) = \mathbb{E}_{\tau \sim \pi_\theta}
  \Big[ \sum_{t=0}^{T} \gamma^{t} \, r_t \Big].
\end{equation}
%
Refer to numbered equations with \verb|\eqref|, as in~\eqref{eq:objective}.
Then describe \method: the key idea, what is new, and why it should work.

%%%%%%%%%%%%%%%%%%%%%%%%%%%%%%%%%%%%%%%%%%%%%%%%%%%%%%%%%%%%%%%%%%%%%%%%

\section{Experiments and Results}
\label{sec:experiments}

Summarize your experimental setup in a few sentences, then report the main results.
For Track 2, please report results separately for each tier. You do not need to redefine
metrics that are already defined as a part of the video to data challenge. It may be
useful to include informative baselines and ablations that support and provide insight
into your claimed contributions.

\begin{table}[t]
  \centering
  \begin{tabular}{lcc}
    \toprule
    \textit{Method} & \textit{Metric 1} & \textit{Metric 2} \\
    \midrule
    Baseline 1      & $0.00 \pm 0.00$ & $0.00 \pm 0.00$ \\
    Baseline 2      & $0.00 \pm 0.00$ & $0.00 \pm 0.00$ \\
    \method (ours)  & $\mathbf{0.00 \pm 0.00}$ & $\mathbf{0.00 \pm 0.00}$ \\
    \bottomrule
  \end{tabular}
  \caption{Main results (mean $\pm$ std over 5 seeds).}
  \label{tab:results}
\end{table}

\begin{figure}[t]
  \centering
  \includegraphics[width=0.6\linewidth]{assets/nvlogo2.png}
  \caption{Overview of \method. Replace with your own figure.}
  \label{fig:overview}
  \Description{Placeholder figure showing the NVIDIA logo.}
\end{figure}

%%% Issue \balance somewhere in the first column of the last page.
\balance

%%%%%%%%%%%%%%%%%%%%%%%%%%%%%%%%%%%%%%%%%%%%%%%%%%%%%%%%%%%%%%%%%%%%%%%%

\section{Conclusion}
\label{sec:conclusion}

Restate the main finding and the most important direction for future work.
Briefly describe any practical implications for your approach at scale.
If a full version of the paper exists, you may point to it here, but this 
is not guaranteed to be reviewed.

%%%%%%%%%%%%%%%%%%%%%%%%%%%%%%%%%%%%%%%%%%%%%%%%%%%%%%%%%%%%%%%%%%%%%%%%

%%% References, acknowledgments and appendices do not count toward
%%% the two-page limit.
\bibliographystyle{plainnat}
\bibliography{references}

%%%%%%%%%%%%%%%%%%%%%%%%%%%%%%%%%%%%%%%%%%%%%%%%%%%%%%%%%%%%%%%%%%%%%%%%

%%% Optional appendix: starts on a new page and does not count toward
%%% the two-page limit. Delete this part if you do not need it.
\clearpage
\appendix

\section{Appendix}
\label{sec:appendix}

\textit{The appendix is optional and is not guaranteed to be read by the
judges. The main text must stand on its own; use this space only for
supporting material such as additional results, implementation details
or hyperparameters.}

\subsection{Additional Results}
\label{sec:appendix-results}

Report supplementary results, additional details, or additional examples here.

\end{document}
